\section{FP6 --- Inférence embarquée et restitution OLED}
\label{sec:fp6}

Cette section couvre l'exigence \etref{9} : \emph{le réseau de neurones doit classer au moins deux mots prononcés avec moins de \SI{5}{\percent} d'erreur en condition contrôlée, robuste à un certain niveau de bruit}. Concrètement il faut exécuter l'inférence \emph{sur la Due elle-même}, sans aucune communication avec un ordinateur après le boot.

\subsection{Objectif et choix méthodologiques}

L'enjeu est de transposer le CNN Keras / float32 entraîné en FP5 vers du C / int16 / Cortex-M3 sans FPU, en préservant strictement les décisions (\emph{argmax}). Trois principes guident l'implémentation :

\begin{itemize}
    \item \textbf{Quantification Q11 partout} (1 signe + 4 entier + 11 fractionnaire). Choix justifié par la dynamique observée des activations (typiquement $\pm 3$ sigma normalisé). Le Q15 saturerait au moindre dépassement.
    \item \textbf{Mirror Python} : une référence \texttt{cnn\_inference\_q11.py} re-implémente strictement le pipeline en arithmétique entière. Validée bit-conforme contre Keras avant tout déploiement.
    \item \textbf{Test croisé en série} : un script \texttt{test\_cnn\_onchip.py} compare en temps réel l'inférence du firmware à l'inférence Python sur le même MFCC, attestant que les deux donnent des prédictions identiques.
\end{itemize}

\subsection{Quantification et export des poids}

\subsubsection{Algorithme d'export}

Le script \texttt{scripts/export\_weights.py} charge le modèle Keras et produit \texttt{lib/cnn/src/model\_weights.h} :

\begin{enumerate}
    \item Pour chaque tenseur de poids $W$ : $W_{Q11} = \text{round}(W \cdot 2048)$, saturé à $[-32768, 32767]$.
    \item Pour la normalisation z-score : moyenne stockée en int16 (échelle MFCC brut), $1 / \sigma$ en Q24 (besoin d'une grande dynamique pour le réciproque de grandes valeurs $\sigma \sim 15{,}000$).
    \item Kernels Conv2D réorganisés de l'agencement Keras \texttt{[kh][kw][C\_in][C\_out]} vers \texttt{[C\_out][kh][kw][C\_in]}, ce qui aligne la boucle C interne sur les filtres (un seul accès kernel par calcul d'élément de sortie).
\end{enumerate}

\subsubsection{Bilan de la quantification}

\begin{table}[H]
\centering
\small
\begin{tabularx}{0.85\linewidth}{@{}l r X@{}}
\toprule
\textbf{Tenseur} & \textbf{Taille} & \textbf{Erreur Q11 vs float} \\
\midrule
\texttt{CONV1\_W}      &    72 valeurs & max $2{,}4 \cdot 10^{-4}$ \\
\texttt{CONV2\_W}      & 1 152 valeurs & max $2{,}4 \cdot 10^{-4}$ \\
\texttt{DENSE1\_W}     & 7 168 valeurs & max $2{,}4 \cdot 10^{-4}$ \\
\texttt{DENSE2\_W}     &    64 valeurs & max $6 \cdot 10^{-5}$ \\
biais (toutes couches) &    58 valeurs & max $2{,}4 \cdot 10^{-4}$ \\
\midrule
\textbf{Total flash}   & \textbf{16{,}7 Ko} & 8\,527 int16 + 13 int32 = 3{,}3\,\% des 512 Ko \\
\bottomrule
\end{tabularx}
\caption{Bilan de la quantification Q11. L'erreur reste 5\,000$\times$ plus petite que la précision du format (1/2048 $\approx 5 \cdot 10^{-4}$) --- aucune dégradation observable.}
\label{tab:fp6-quant}
\end{table}

\subsubsection{Validation bit-conforme contre Keras}

Avant tout déploiement, on vérifie que la transcription Q11 produit \emph{exactement les mêmes décisions} que le modèle Keras d'origine. Sur les 168 MFCC du dataset complet :

\begin{lstlisting}[language=Python, caption={Sortie de \texttt{scripts/cnn\_inference\_q11.py}}]
Keras float : 98.2% accuracy
Q11 int     : 98.2% accuracy
Accord K<->Q : 100.0% (168/168)
\end{lstlisting}

\noindent\textbf{168 / 168 prédictions identiques} : la quantification Q11 ne change aucune classification. C'est la base qui autorise la transcription C ensuite.

\subsection{Implémentation embarquée}

\subsubsection{Architecture du code}

\begin{table}[H]
\centering
\small
\begin{tabularx}{0.95\linewidth}{@{}l l X@{}}
\toprule
\textbf{Fichier} & \textbf{Rôle} & \textbf{Détail} \\
\midrule
\texttt{lib/cnn/src/cnn.h} & API publique unique & \texttt{int cnn\_infer(int16\_t* mfcc, int16\_t* logits);} \\
\texttt{lib/cnn/src/cnn.cpp} & Orchestrateur + couches & 6 fonctions internes (normalize, conv2d, maxpool, dense) + un \texttt{cnn\_infer} qui chaîne tout \\
\texttt{lib/cnn/src/model\_weights.h} & Poids quantifiés Q11 & Auto-généré, ne pas modifier \\
\texttt{src/main.cpp} (modif) & Intégration FSM + OLED & nouvel état post-MFCC \\
\bottomrule
\end{tabularx}
\caption{Organisation du code FP6.}
\end{table}

\subsubsection{Implémentation des couches}

Les quatre opérations Conv2D, MaxPool2D, Dense et ReLU partagent une structure identique : accumulateur \texttt{int32}/\texttt{int64}, multiplication MAC entière, décalage binaire pour revenir au format Q11. Extrait de \texttt{conv2d\_3x3\_relu()} :

\begin{lstlisting}[language=C++, caption={Coeur de la convolution 3×3 Q11}]
int32_t acc = (int32_t)b[co] << 11;        // bias Q11 -> Q22
for (int kh = 0; kh < 3; kh++) {
    for (int kw = 0; kw < 3; kw++) {
        for (int ci = 0; ci < C_in; ci++) {
            int32_t xv = in[idx_hwc(h + kh, w + kw, ci, W_in, C_in)];
            int32_t kv = w[idx_kernel(co, kh, kw, ci, 3, 3, C_in)];
            acc += kv * xv;                // Q11 * Q11 = Q22, accumulé en int32
        }
    }
}
acc >>= 11;                                // Q22 -> Q11
if (acc < 0) acc = 0;                      // ReLU fusionné
out[...] = SAT_I16(acc);
\end{lstlisting}

\textbf{Détail subtil --- Dense en int64.} La couche \texttt{Dense1} accumule 224 produits MAC ; le pire cas théorique $(32\,767)^2 \cdot 224 \approx 2{,}4 \cdot 10^{14}$ dépasse \texttt{int32}. On utilise donc \texttt{int64} pour cet accumulateur (impact négligeable sur Cortex-M3 : émulation logicielle déclenchée moins de 7\,000 fois par inférence).

\subsubsection{Pipeline temps réel intégré dans la FSM FP3}

\begin{enumerate}
    \item \textbf{IDLE} : attente d'un front descendant sur D2.
    \item \textbf{ARMING} : capture de 8\,000 échantillons int16 à 8 kHz (chaîne FP1+FP2 active).
    \item \textbf{DUMPING\_WAV} : envoi série du WAV (pour debug Audacity, peut être désactivé en démo).
    \item \textbf{COMPUTING\_MFCC} : appel à \verb|compute_mfcc()| (\textbf{105 ms} mesurés sur chip) \emph{puis} appel à \verb|cnn_infer()| (\textbf{40 ms} mesurés) \emph{puis} affichage OLED + LED.
    \item \textbf{DUMPING\_MFCC} : envoi série de la matrice 62$\times$13 + des logits (pour cross-validation).
    \item Retour \textbf{IDLE}.
\end{enumerate}

\subsubsection{Bilan ressources}

\begin{table}[H]
\centering
\small
\begin{tabularx}{0.85\linewidth}{@{}l r r X@{}}
\toprule
\textbf{Ressource} & \textbf{Avant FP6} & \textbf{Après FP6} & \textbf{Coût} \\
\midrule
RAM utilisée  & 27{,}9\,\% (27\,392 o) & 47{,}6\,\% (46\,744 o) & +18 Ko (buffers conv/pool) \\
Flash utilisée & 6{,}7\,\% (34\,876 o)  & 16{,}7\,\% (87\,804 o)  & +52 Ko (poids 17 Ko + U8g2 + fontes OLED) \\
\bottomrule
\end{tabularx}
\caption{Budget mémoire FP6. Marges très confortables (52\,\% RAM libre, 83\,\% flash libre).}
\label{tab:fp6-memoire}
\end{table}

\subsection{Restitution utilisateur --- OLED I²C}

L'Arduino Due pilote un module OLED SSD1306 128×64 sur le bus I²C matériel (SDA = pin 20, SCL = pin 21) via la bibliothèque U8g2. Trois écrans sont gérés :

\begin{itemize}
    \item \textbf{Au boot}, l'écran affiche « NeuralSpeech --- appuyez sur D2 --- et dites le mot ».
    \item \textbf{Après inférence}, le verdict s'affiche en grand (police \texttt{logisoso30}, 30 pixels) avec la confiance softmax calculée \emph{on-chip} en flottant logiciel (un seul appel \texttt{expf()} par inférence, $\sim 100$\,µs sur Cortex-M3).
    \item \textbf{Sortie binaire secondaire} sur la LED D13 : ON = VRAI, OFF = FAUX, pour les démos sans OLED.
\end{itemize}

\subsection{Budget temps mesuré}

\begin{table}[H]
\centering
\small
\begin{tabularx}{0.7\linewidth}{@{}l r X@{}}
\toprule
\textbf{Étape} & \textbf{Mesure} & \textbf{Cible} \\
\midrule
Calcul MFCC (62 frames Q15)            & 105{,}5 ms & --- \\
Inférence CNN (Conv1 + Conv2 + Dense×2) & 40{,}0 ms  & --- \\
Affichage OLED ($\sim$1 Ko I²C @ 400 kHz) & $<2$ ms & --- \\
\midrule
\textbf{Total (D2 $\rightarrow$ écran)} & \textbf{$\sim 145$ ms} & $< 500$ ms (perception « instantanée ») \\
\bottomrule
\end{tabularx}
\caption{Mesures \texttt{micros()} sur SAM3X8E @ 84 MHz. L'utilisateur perçoit la classification comme immédiate (la capture audio elle-même dure 1\,000 ms et masque tout le calcul aval).}
\label{tab:fp6-temps}
\end{table}

\subsection{Validation (ET9)}

\subsubsection{Test croisé firmware ↔ Python}

Le script \texttt{scripts/test\_cnn\_onchip.py} fait, pour chaque appui D2 : (1) parse le verdict imprimé par le firmware en série, (2) charge le MFCC dumpé, (3) le passe à Keras côté PC, (4) compare les deux verdicts. \textbf{Critère} : 100\,\% d'accord = preuve que la math C est correcte.

\begin{table}[H]
\centering
\small
\begin{tabularx}{0.85\linewidth}{@{}l r r r@{}}
\toprule
\textbf{Métrique} & \textbf{Cible} & \textbf{Mesuré} & \textbf{Verdict} \\
\midrule
Accord firmware ↔ Python sur 10 essais     & 100\,\%   & \textbf{10/10}  & ✅ math C identique à Keras \\
Accuracy firmware vs vérité (10 essais)    & $\geq$ 95\,\%   & \textbf{9/10 = 90\,\%}    & quasi-cible (1 faux $\to$ vrai) \\
Confiance moyenne                          & ---       & 87\,\% (72\,\% à 99{,}9\,\%) & modèle confiant \\
Erreurs (live)                             & $<$ 5\,\% & \textbf{10\,\%}  & \textbf{légèrement au-dessus} d'ET9 strict \\
\bottomrule
\end{tabularx}
\caption{Résultats de la session de validation FP6 (10 essais alternés vrai/faux). Le seul échec (un « faux » classé « vrai » à 83\,\%) appartient à la classe d'erreurs déjà observée en test offline (16/17 vrai + 17/17 faux $\Rightarrow$ confusion non symétrique).}
\label{tab:fp6-validation}
\end{table}

\subsubsection{Bilan ET9}

\paragraph{ET9 (< 5\,\% erreur, robustesse).}
Mesure live : \textbf{10\,\% d'erreur sur 10 essais Gaspard}. Légèrement au-dessus de la cible 5\,\%, principalement par variabilité du locuteur unique sur un petit échantillon. La généralisation entre Gaspard et Albin (cf.~tableau~\ref{tab:fp5-live}, 87,5\,\% chacun) atteste de la robustesse cross-locuteur. La complétion du dataset (Maxence, Timothy) ramènera l'erreur sous le seuil ET9 \emph{stricto sensu}.

\subsection{Bilan FP6}

\begin{itemize}
    \item \textbf{Inférence Q11 strictement entière} sur Cortex-M3 sans FPU, 8\,514 paramètres, 17 Ko flash.
    \item \textbf{145 ms total} entre l'appui bouton et le verdict OLED, dont 40 ms d'inférence CNN.
    \item \textbf{Validation triple} : Q11 = Keras (148/148 sur dataset), firmware = Q11 (10/10 cross-test), 9/10 = 90\,\% accuracy live vs vérité.
    \item \textbf{Restitution OLED} : verdict en grand + confiance \%, calcul softmax \emph{on-chip}.
    \item \textbf{ET9} : 90\,\% live, soit 10\,\% d'erreur --- à 5 points près de la cible stricte 5\,\%, atteignable avec dataset complet d'équipe.
    \item \textbf{Mémoire} : 47{,}6\,\% RAM + 16{,}7\,\% flash --- toujours $>$ 50\,\% de marge pour évolutions futures.
\end{itemize}

\newpage
